\section{Simplicial sets}\label{sec3}

So far, we have learned about geometric simplices, their cofaces and codegeneracies, and how the simplex category $\BD$ encodes this information combinatorially. As promised, we now explore simplicial sets as a more powerful way to encode combinatorial data about various structures. Like the simplex category, simplicial sets are closely related to topological spaces and the geometric simplicies because they admit a \emph{geometric realization}, which is a process for constructing topological spaces out of simplicial sets. The use of simplicial sets is ubiquetous throughout algebraic topology and homotopy theory, and we will see some applications in \cref{sec4}.

In this section, we will have to consider universal constructions such as quotients and coproducts in the category $\Top$. For the topologically uninitiated reader, we make the following remark:

\begin{remark}\label{3.0.1}
    There is a forgetful functor $U:\Top\to\Set$ which admits both a left \emph{and} a right adjoint. In particular, $U$ preserves all limits and colimits. Since $U$ is a forgetful functor, it does not change the underlying set of a topological space, and thus all universal constructions in $\Top$ have the same underlying set as those in $\Set$.
\end{remark}


\subsection{Definition}\label{sec3.1}
The definition of simplicial sets is fairly innocuous, but the level of detail that our problems demand we work in will reveal the definition to be suprisingly complex. Most experts find simplicial sets very simple, but students often take significant time to fully process the definition and to come to terms with it. We develop many examples and exercises throughout this section to improve the reader's understanding. We highly recommend to work through all the exercises.

\begin{definition}\label{3.1.1}
    A simplicial set $X_{\bullet}$ is a contravariant functor $X : \BD\opp\to\Set$. We generally denote the image $X[n]$ on an object by the subscript $X_n$, hence the bullet notation. We will generally not denote the action on morphisms by subscripts, and instead appeal to \cref{3.1.2}.
\end{definition}

\begin{notation}\label{3.1.2}
By \cref{2.3.4}, we need only consider the action of a simplicial set $X$ on cofaces and codegeneracies. We denote $X(d^i)$ by $d_i$, and call it a \emph{face map}. Similarly, we denote $X(s^i)$ by $s_i$, and call it a \emph{degeneracy map}. The term ``simplicial operators'' refers to faces, degeneracies, and in general any map in the image of a simplical set $X$.
\end{notation}

\begin{remark}\label{3.1.3}
    Note that since $X : \BD\opp\to\Set$ is contravariant, face maps and degeneracy maps go in the opposite direction compared to the coface maps and codegeneracy maps in $\BD$. More precisely, $d_i : X_n\to X_{n-1}$ and $s_i : X_n\to X_{n+1}$.
\end{remark}

\begin{intuition}\label{3.1.4}
    Note that a simplicial set is really just a function which takes each number $[n]$ and assigns it to a set $X_n$. Of course these sets have to satisfy some properties encoded by the faces and degeneracies, which is discussed in \cref{3.1.7}. The correct way to think about a simplicial set is thus as a collection of sets $X_n$ for each $n$, which comes imbued with some extra structure given by the simplicial operators.
\end{intuition}

\begin{remark}\label{3.1.5}
    Since simplicial sets are just functors $\BD\opp\to\Set$, we can construct a \emph{category of simplicial sets} by simply considering the functor category $\Fun(\BD\opp,\Set)=\Set^{\BD\opp}$. We will generally denote this category as $\sSet$. Note that morphisms in this category are given by natural transformations, so we define morphisms $\alpha : X_{\bullet} \to Y_{\bullet}$ of simplicial sets to simply be natural transformations from $X$ to $Y$. By \cref{2.3.4}, we only need to consider faces and degeneracies. We thus have that a natural transformation $\alpha : X_{\bullet} \to Y_{\bullet}$ is a collection of arrows $\alpha_n : X_n \to Y_n$ for each $[n]\in\BD$ such that the naturality squares
    \[
    \begin{tikzcd}
        X_n \ar["\alpha_n", r]\ar["d_i"', d] & Y_n \ar["d_i", d] & X_n \ar["\alpha_n", r] \ar["s_i"', d] & Y_n \ar["s_i", d]\\
        X_{n-1} \ar["\alpha_{n-1}"', r] & Y_{n-1} & X_{n+1} \ar["\alpha_{n+1}"', r] & Y_{n+1}
    \end{tikzcd}
    \]
    commute. In other words, a morphism $\alpha : X_{\bullet} \to Y_{\bullet}$ is a collection of maps $\alpha_n$ for each $n$ that commute with faces and degeneracies.
\end{remark}

\begin{terminology}\label{3.1.6}
    Given a simplicial set $X_{\bullet}$, we call an element $x\in X_n$ an $n$-\emph{simplex}. As in \cref{2.1.2}, the plural of $n$-simplex is ``$n$-simplices''. We justify this terminology in \cref{3.2.6}. An $n$-simplex $x\in X_n$ is \emph{degenerate} if it is in the image of $\emph{any}$ degeneracy map $s_i : X_{n+1}\to X_n$. Given an $n$-simplex $x\in X_n$, the $i$th \emph{face} of $x$ is the image of the face map $d_i(x)$.
\end{terminology}

\begin{remark}[Simplicial identities]\label{3.1.7}
    We pause to investigate faces and degeneracies in simplicial sets. We have seen in \cref{2.3.4} and \cref{2.3.6}, as summarized in \cref{2.3.8}, that the cofaces and codegeneracies, when subject to the cosimplicial operators, generate all morphisms in $\BD$. Similarly, any specified collection of faces $d_i : X_n \to X_{n-1}$ and degeneracies $s_i : X_n\to X_{n+1}$ will generate all simplicial operators for a simplicial set $X$, when subject to the following \emph{simplicial identities}:
    \[
    \begin{cases}
        d_id_j = d_{j-1}d_i,& i<j\\
        d_is_j = s_{j-1}d_i, & i<j\\
        d_js_j=d_{j+1}s_j=1&\\
        d_is_j=s_jd_{i-1},&i>j+1\\
        s_is_j=s_{j+1}s_i,&i\leq j
    \end{cases}.
    \]
    Note that these are simply the duals of \cref{2.3.6}, and indeed they follow by contravariance and functoriality of $X$. Thus, to define a simplicial set $X$, it suffices to do the following three things:
    \begin{enumerate}
        \item Define a set $X_n$ for each $[n]\in\BD$;
        \item For each $i\in[n]$, define a map $d_i : X_n\to X_{n-1}$ and a map $s_i : X_n \to X_{n+1}$;
        \item Show the maps of (2) satisfy the simplicial identities above.
    \end{enumerate}
    These are precisely the qualities you obtain by forcing $X$ to be functorial.
\end{remark}

\begin{exercise}\label{3.1.8}
    We write the representable functor $\BD(-,[n])$ as $\Delta^n$; the same notation for the standard $n$-simplex. We then write the Yoneda embedding $\BD\to\sSet$ as $\Delta^\bullet$. The reasoning for this notation will become clear in \cref{3.2.7}, but for now, complete the following exercises.
    \begin{enumerate}
        \item Show that the $n$-simplicies $\Delta^0_n=*$ for all $n$, where $*$ denotes the singleton. In particular, show that $\Delta^0$ is terminal in $\sSet$.
        \item We call the simplicial set $\Delta^n$ the standard $n$-simplicial set, or sometimes the standard $n$-simplex, although we think the latter is stupid because it conflates notation even more with topological $n$-simplices. Yoneda's lemma tells us quite a lot about the standard simplicial sets for free. What can you deduce about these objects using only the Yoneda lemma, and basic category theory? Make sure part of your answer includes a description of faces and degeneracies for $\Delta^n$.
    \end{enumerate}
\end{exercise}

\begin{example}\label{3.1.9}
    Any set $X$ can be regarded as a simplicial set by taking each $X_n=X$, and by taking the faces and degeneracies to all be the identity map $1_X$. This definition manifeslty satisfies the simplicial identities as stated in \cref{3.1.7}, indeed making $X$ a simplicial set. We call such a simplicial set \emph{discrete}, and this is the canonical way to view any set as a simplicial set. We will sometimes use sets in places where simplicial sets are supposed to be used. In those cases, we are implicitly considering the set $X$ as a discrete simplicial set.
\end{example}

\begin{definition}\label{3.1.10}
	Let $X_\bullet$ be a simplicial set, and recall \cref{3.1.6}.
    \begin{enumerate}
        \item The simplicial set $X_\bullet$ is said to be $n$-\emph{skeletal} if for $m>n$, all $m$-simplices of $X_\bullet$ are degenerate.
        \item The $n$-truncated simplex category $\BD_{\leq n}$ is the full subcategory of $\BD$ with objects $[m]$ for $m\leq n$.
        \item The $n$-truncation $X_{\leq n}$ of a simplicial set $X_\bullet$ is the restriction of the simplicial set to the $n$-truncated simplex category $X|_{\BD_{\leq n}}$, which is intuitively understood as throwing away all dimensions above $n$.
        \item The $n$-skeleton $\mathrm{sk}_n(X)$ of a simplicial set $X_\bullet$ is obtained by removing all nondegenerate $m$-simplices of $X$ for $m>n$. This is often phrased as taking the $n$-truncation and replacing higher dimensions $m>n$ with freely generated degenerate simplices.
    \end{enumerate}
\end{definition}

\subsection{Geometric realization}\label{sec3.2}
We were able to view objects in $\BD$ as geometric objects using the functor $\Delta^{\bullet}  : \BD\to\Top$. We can in fact achieve even better results with the \emph{geometric realization functor} $|-| : \sSet\to\Top$. What follows in \cref{3.2.1} is an extremely succinct way to summarize the geometric realization of a simplicial set using the language of \emph{coends}.

Although it is simple, it is not particularly descriptive. Thus the reader who is unfamiliar with coends should not be concerned, as we will provide a much better description in \cref{3.2.2} and \cref{3.2.3}. We found it suprisingly difficult to find a good resource on geometric realization; most sources expect us to already know it, yet we cannot find where we are supposed to know it from. We thank Emily Riehl for her presentation in \cite{sset}, and we thank Greg Friedman for his presentation in \cite{fried}.

\begin{definition}\label{3.2.1}
    Let $X_\bullet$ be a simplicial set, and let $\Delta^n$ denote the standard $n$-simplex as defined in \cref{2.1.1}. The geometric realization $|X_\bullet|$ of $X_\bullet$ is the coend
    \[
    \int^{[n]\in\BD} X_n \times \Delta^n,
    \]
    where $X_n$ is given the discrete topology. This is equivalent to the definition via functor tensor products, since the usual tensor is the levelwise copower, and there is a homeomorphism $X_n\cdot\Delta^n\cong X_n\times\Delta^n$ when $X_n$ is given the discrete topology.
\end{definition}

\begin{construction}\label{3.2.2}
    We now detail the construction of the geometric realization of a simplicial set. The experienced reader should be able to recognize this as the coend given in \cref{3.2.1}, computed via coproducts and coequalizers. Let $X_\bullet$ be a simplicial set, and recall that the coproduct in $\Top$ is the disjoint union. The geometric realization $|X_\bullet|$ is computed by following these steps:
    \begin{enumerate}
        \item Construct the coproduct
        \[
        \coprod_{[n]\in\BD} X_n\times\Delta^n,
        \]
        where the set $X_n$ is given the \emph{discrete topology}, meaning all subsets $U\subseteq X_n$ are open.
        \item We then impose an equivalence relation on the space, which we plan to quotient by. Let $\alpha : [m]\to[n]$ be a morphism in the simplex category $\BD$. Denote by $\alpha^*$ the morphism $X(\alpha)$, and denote by $\alpha_*$ the continuous map $\Delta^\bullet(\alpha)$. We describe in more detail how to think of these morphisms in \cref{3.2.3}. We quotient by the following relation: for any $x\in X_n$ and any $p\in \Delta^m$,
        \[
        (\alpha^*(x),p)\sim (x,\alpha_*(p)).
        \]
        Note that this quotient only affects the parts of the coproduct indexed by $[n]$ and $[m]$, so we intuitively think of it as only quotienting those parts. We then impose this quotient \emph{for all} morphisms $\alpha$, which gives the geometric realization $|X_\bullet|$.
    \end{enumerate}
\end{construction}

\begin{remark}\label{3.2.3}
    Of course by \cref{2.3.4}, we only need to consider cofaces and codegeneracies, not \emph{all} morphisms $\alpha$ in $\BD$. This simplifies the construction substantially, and we mainly described \cref{3.2.2} in terms of all morphisms for pedagogical purposes; in particular to emphasize that cofaces and codegeneracies subsume all other morphisms in $\BD$. We continue with the understanding that we need only work with cofaces and codegeneracies.
    \begin{itemize}
        \item (faces) Recall that the action of $X$ on the coface map $d^i : [n-1]\to[n]$ is given by the face map $d_i : X_n\to X_{n-1}$. Recall by \cref{2.4.1} and \cref{2.4.4} that the corresponding coface map $d^i:\Delta^{n-1}\to\Delta^{n}$ inserts $\Delta^{n-1}$ at the $i$th face of $\Delta^n$. What we want to do is think of each $n$-simplex $x\in X_n$ as identifying a geometric $n$-simplex via $\{x\}\times \Delta^n\subseteq X_n\times\Delta^n$. The product space can thus be thought of as an $X_n$-indexed (disjoint) collection of geometric $n$-simplices. The quotient process is thus interpreted as follows:
        \begin{enumerate}
            \item Choose an $n$-simplex $x\in X_n$, and find the corresponding geometric $n$-simplex $\Delta^n$.
            \item Find the $i$th face $d_i(x)\in X_{n-1}$ of $x$, and find the corresponding geometric $(n-1)$-simplex $\Delta^{n-1}$.
            \item Glue the $(n-1)$-simplex $\Delta^{n-1}$ corresponding to $d_i(x)$ to the $i$th face of the geometric $n$-simplex $\Delta^n$ corresponding to $x$ by gluing along the coface map $d^i : \Delta^{n-1}\to\Delta^n$.
        \end{enumerate}
        \item (degeneracies) Recall that the action of $X$ on the codegeneracy $s^i : [n+1]\to[n]$ is given by the degeneracy map $s_i : X_n\to X_{n+1}$. Recall by \cref{2.4.1} and \cref{2.4.4} that the corresponding codegeneracy map $s^i : \Delta^{n+1}\to\Delta^n$ identifies the $i$th and $(i+1)$th vertices of $\Delta^{n+1}$, in effect ``squishing'' the simplex down to a degenerate $n$-simplex. Using the same logic as we did for faces, we note that each $n$-simplex $x\in X_n$ identifies a geometric $n$-simplex $\{x\}\times\Delta^n$, so the quotienting process is interpreted as follows:
        \begin{enumerate}
            \item Choose an $n$-simplex $x\in X_n$, and find the corresponding geometric $n$-simplex $\Delta^n$.
            \item Find a degenerate $(n+1)$-simplex $s_i(x)\in X_{n+1}$, and find the corresponding geometric $(n+1)$-simplex $\Delta^{n+1}$.
            \item Collapse the $(n+1)$-simplex $\Delta^{n+1}$ corresponding to $s_i(x)$ to the $n$-simplex $\Delta^n$ corresponding to $x$ along the $i$th codegeneracy $s^i$, identifying $\Delta^{n+1}$ with a degenerate simplex.
        \end{enumerate}
        We can summarize this process as taking all the geometric simplices corresponding to degenerate simplices, and turning them into degenerate simplices by composing along the corresponding codegeneracy.
    \end{itemize}
    The resulting information we gather from this analysis is twofold. Degenerate simplices can be thought of as \emph{redundant information}, since the corresponding geometric simplex collapses down to a lower dimensional simplex. Faces of a simplex can be thought of as geometric faces of the corresponding geometric simplex, again by the process described above.
\end{remark}

\begin{intuition}\label{3.2.4}
	The geometric realization, at first glance, is a fairly contrived and arbitrary construction. However, now that we've explored the construction in \cref{3.2.2} and \cref{3.2.3}, we can provide some decent intuition for what simplicial sets really are.

	Recall that in \cref{sec2.4}, we obtained an understanding of how $\boldsymbol{\Delta} $ modeled simplicial sets and their relationships between each other in a combinatorial way by identifying each element with an $n$-simplex, and identifying each coface and codegeneracy with the corresponding geometric coface and codegeneracy. The category $\sSet $ is more powerful because it encodes \emph{more} combinatorial information that $\boldsymbol{\Delta} $. In \cref{3.2.7}, we will show the geometric realization of the standard $n$-simplicial sets are the standard $n$-simplices, and we will see even later that there is a notion of coface and codegeneracy between the standard $n$-simplicial sets, which correspond precisely to the notions of the same name between the standard geometric $n$-simplices. In this way, the objects $\Delta^n$ and the cofaces and codegeneracies between them encode the same information as $\boldsymbol{\Delta} $, since geometric realization gives us the same objects and the same morphisms.

	However, $\sSet $ is a much larger category than just the simplicial sets given by the Yoneda embedding. Other simplicial sets can have simpler structure, or more complicated structure, allowing us to build much more interesting spaces. Each simplicial set $X_{\bullet} $ encodes combinatorial information via its faces and degeneracies, and this information tells the geometric realization how to construct a geometric object. It starts by building geometric $n$-simplices for each $n$-simplex of $X_{\bullet} $, and then it glues them together based on the combinatorial information $X_{\bullet} $ contains. The vast amount of freedom we have when defining a simplicial set allows us to construct a large variety of different spaces. For example, \cref{3.3.5} constructs a space known as the ``classifying space'' for the very simple group $\mathbb{Z} /2\mathbb{Z} =\mathbb{Z} _2$, which is the geometric realization of a related simplicial set called the \emph{nerve}. We find that the classifying space is the infinite real projective space $\mathbb{R} P^\infty $. Even though the group $\mathbb{Z} /2\mathbb{Z} $ is a very simple group, we are still able to construct an extremely complex space using the information it encodes. And we can also construct simple spaces with simplicial sets: in \cref{3.3.9} you will construct the discrete topological space on any set $X$ in two different ways.

	Thus, simplicial sets are ways to encode extremely complex data in a simple combinatorial fasion. It's a set of rules for gluing and collapsing geometric simplices, as described in \cref{3.2.2}, which can produce complex and important spaces. As we saw in \cref{sec2.4}, this can simply quite a number of proofs, but simplicial objects have applications far beyond just this. In this paper, we will explore only a small number of the applications of simplicial sets, mainly with a focus on homotopy theory and category theory. However, applications of simplicial sets have become somewhat ubiquitous in algebraic topology, category theory, and fields beyond these in the past few years due to their \emph{relatively} simple nature, compared to the complex structure they encode.

    The reader at this point is likely hoping for some examples of geometric realization computations. However, the geometric realization can be fairly lengthy to compute for nontrivial cases, and many computations require theory that will not be developed until later, if at all. For this reason, most of our examples of computations will center around developing some of this theory, and will be sprinkled throughout the text. Our first computation will be \cref{3.2.7}, where we compute the geometric realization of the standard $n$-simplicial sets $\Delta^n$. For the reader who immediately wants to see an example of a more advanced worked geometric realization, we recommend skipping to \cref{sec3.3}, and reading \cref{3.3.5}.
\end{intuition}

\begin{exercise}\label{3.2.5}
    We claim in \cref{3.2.3} that by \cref{2.3.4}, we need only consider the faces and degeneracies when computing the geometric realization of a simplicial set. Make this notion precise by proving that if we construct the relations of \cref{3.2.2} (2) for only the faces and degeneracies, then we already necessarily have the relation $(\alpha^*(x),p)\sim(x,\alpha_*(p))$ for any arrow $\alpha$ in $\BD$, with notation as in \cref{3.2.2}.
\end{exercise}

\begin{remark}\label{3.2.6}
    Throughout \cref{3.2.3} and \cref{3.2.5}, we continuously thought of each product $X_n\times\Delta^n$ not quite as a product space, but instead as an $X_n$-indexed collection of geometric $n$-simplices. The reader may be wondering if this is equivalent to the copower $X_n\cdot\Delta^n$, especially since these constructions are isomorphic as sets. Indeed, when $X_n$ is given the discrete topology, the copower is homeomorphic to the product. This is not difficult to prove, and we encourage the reader to work out the proof if they are not convinced. We typically prefer product notation since it is more simple to write $(x,-)$ to denote that $x$ is the indexing $n$-simplex in $X_n$, which simplifies the presentation of the relation we quotient by in \cref{3.2.2}. However this is largely notational, and we still think of it more as a copower.
\end{remark}

\begin{proposition}\label{3.2.7}
    The geometric realization of the standard $n$-simplicial set $\Delta^n$, as defined in \cref{3.1.9}, is the standard $n$-simplex $\Delta^n$.
\end{proposition}
\begin{proof}
    Due to our overloaded notation, we will need to adopt some new notational conventions for the duration of this proof in order to make it readable. We will denote the standard $n$-simplicial set as $\Delta^n$, just like before, and will denote the standard geometric $n$-simplex as $\Delta_n$, with a subscript. Any time both a superscript \emph{and} a subscript show up of the form $\Delta^n_k$, this denotes the $k$-simplices of $\Delta^n$. Furthermore, the functor $\Delta_\bullet$ will denote the canonical embedding $\BD\to\Top$, and $\Delta^\bullet$ will be the Yoneda embedding $\BD\to\Set^{\BD\opp}$.

    An element of $\left|\Delta^n\right|$ is an equivalence class of a pair $(f,p)$, where $f\in\Delta_k^n$ and $p\in\Delta_k$. By the defintion of $\Delta^n$ as given in \cref{3.1.8}, $k$-simplices are morphisms $f : [k]\to[n]$. Note that by \cref{2.4.2}, $f$ induces a continuous map $\Delta_\bullet(f) : \Delta_k\to\Delta_n$. Since $p\in\Delta_k$, we can evaluate this morphism at $p$. Thus we propose the map $\varphi : \left|\Delta^n\right| \to\Delta_n$ given by $[(f,p)]\mapsto \Delta_\bullet(f)(p)$. It is not obvious that this map is well-defined, so before we can show it is a homeomorphism, we must show that any choice of representative of the equivalence class evaluates to the same point.

    We do this by showing $\varphi $ is given by taking a function $\psi $ which respects $\sim $, and descending along the quotient. For all $0\leq k\leq n$, and any $f\in\Delta^n_k$, $p\in\Delta_k$, and $\alpha : [h]\to[k]$ in $\BD$, define the map $\psi$ by
    \[\begin{tikzcd}[row sep=0em]
	{\displaystyle{\coprod_{[k]\in\BD}}\Delta^n_k\times\Delta_k} & {\Delta_n} \\
	{(f,p)} & {\Delta_\bullet(f)(p).}
	\arrow["\psi", from=1-1, to=1-2]
	\arrow[maps to, from=2-1, to=2-2]
    \end{tikzcd}\]
     Let $p'\in \Delta _h$. By \cref{3.2.2}, the equivalence relation we quotient by identifies points $(\alpha^*(f),p')\sim(f,\alpha_*(p'))$, where $\alpha^*$ is the simplicial operator $\Delta^n_k\to\Delta^n_h$ induced by $\Delta^n$, and $\alpha_* : \Delta_h\to\Delta_k$ is given by $\Delta_\bullet(\alpha)$. Recall that since $\Delta^n$ is given by the Yoneda embedding, the morphism $\alpha^*$ is simply given by precomposition, so we obtain that the relation reads as
    \[
    (f\circ \alpha,p')\sim(f,\Delta_\bullet(\alpha)(p')).
    \]
    Applying $\psi$ to both sides of this quotient gives, by functoriality of $\Delta_\bullet$,
    \[
    	\psi (f\circ \alpha ,p') = \Delta_{\bullet} (f\circ \alpha )(p'),
    \]
    \begin{align*}
	    \psi (f, \Delta _{\bullet} (\alpha )(p'))&=\Delta _{\bullet} (f)(\Delta _{\bullet} (\alpha )(p'))\\
						     &=(\Delta _{\bullet} (f)\circ \Delta _{\bullet} (\alpha ))(p')\\
						     &=\Delta _{\bullet} (f\circ \alpha )(p')
    .\end{align*}
    Since these are equal, we have that $\psi$ identifies elements under the equivalence relation $\sim$, and thus by univeraslity of the quotient, there exists a unique dashed arrow
    \[\begin{tikzcd}
	{\displaystyle{\coprod_{[k]\in\BD}}\Delta^n_k\times\Delta_k} & {\Delta_n} \\
	{\displaystyle{\left|\Delta^n\right|\cong\left(\coprod_{[k]\in\BD}\Delta^n_k\times\Delta_k\right)/\tildesim}}
	\arrow["\psi", from=1-1, to=1-2]
	\arrow[from=1-1, to=2-1]
	\arrow[dashed, from=2-1, to=1-2]
    \end{tikzcd}\]
    rendering the diagram commutative. It's important to note that we have not show $\psi $ is continuous yet, so this diagram only exists in $\Set $. By our extensive knowledge of universals in $\Set$, we obtain an explicit description of the dashed arrow as $[(f,p)]\mapsto \psi(f,p)=\Delta_\bullet(f)(p)$, which is precisely the definition of $\varphi$. Thus $\varphi$ arises by universality, and is therefore well-defined.

    Now we show $\varphi$ is continuous. Since $\varphi$ arises by universality, and since by \cref{3.0.1} all universals in $\Top$ admit the same point-set definition as in $\Set$, it suffices to show $\psi$ is continuous, since the quotient morphism is continuous by definition of the quotient topology. This will promote the universal problem to $\Top$, rendering $\varphi$ continuous, as required. We start by considering the restriction of $\psi$ to only one component of the product $\Delta_k^n\times\Delta_k$. We can decompose the restriction of $\psi$ as the composite
    \[\begin{tikzcd}
	{\Delta_k^n\times\Delta_k} & {\Top(\Delta_k,\Delta_n)\times\Delta_k} & {\Delta_n.}
	\arrow["{\Delta_\bullet\times1}", from=1-1, to=1-2]
	\arrow["{\mathrm{eval}}", from=1-2, to=1-3]
    \end{tikzcd}\]
    If we give $\Top(\Delta_k,\Delta_n)$ the compact-open topology, then it is well known that the evaluation map $\mathrm{eval}(\Delta_\bullet(f),p)=\Delta_\bullet(f)(p)$ is continuous whenever the domain (in this case, $\Delta_k$) is locally compact Hausdorff, meaning for every $p\in \Delta _k$ there is an open neighborhood $U$ and a compact neighborhood $K$ with $p\in U \subseteq K$. Indeed, all Euclidean spaces $\mathbb{R}^n$ are locally compact Hausdorff by Heine-Borel, and all open or closed subsets of locally compact Hausdorff spaces are locally compact Hausdorff. Since $n$-simplices are clearly closed as subspaces of $\mathbb{R}^{n+1}$, this follows. It suffices to show $\Delta_\bullet\times 1$ is continuous. Note that $\Delta_\bullet\times 1$ arises as the unique dashed arrow in the diagram
    \[\begin{tikzcd}
	{\Delta_k^n} & {\Top(\Delta_k,\Delta_n)} \\
	{\Delta_k^n\times\Delta_k} & {\Top(\Delta_k,\Delta_n)\times\Delta_k} \\
	{\Delta_k} & {\Delta_k}
	\arrow["{\Delta_\bullet}", from=1-1, to=1-2]
	\arrow[from=2-1, to=1-1]
	\arrow["{\Delta_\bullet\times1}", dashed, from=2-1, to=2-2]
	\arrow[from=2-1, to=3-1]
	\arrow[from=2-2, to=1-2]
	\arrow[from=2-2, to=3-2]
	\arrow["1"', from=3-1, to=3-2]
    \end{tikzcd}\]
    induced by universality of products in $\Top$. Since the identity $1$ is continuous, it suffices to show $\Delta_\bullet$ is continuous. Recall that all functions whose domain is a discrete topological space are continuous, since the preimage of any open set is a set, and any set is open in a discrete space. Recall in \cref{3.2.1,3.2.2} that we grant the set of $k$-simplices $\Delta_k^n$ the discrete topology, and thus $\Delta_\bullet$ is continuous. We thus have that the restriction of $\psi$ is continuous for any $k$.

    Since every restriction $\psi|_k$ of $\psi$ along the $k$-indexed component of the coproduct is continuous, universality of coproduct assembles the arrows $\psi|_k$ into a unique dashed arrow as indicated
    \[\begin{tikzcd}
	{\Delta_k^n\times\Delta_k} \\
	{\displaystyle{\coprod_{[k]\in\BD}\Delta_k^n\times\Delta_k}} & {\Delta_n}
	\arrow[from=1-1, to=2-1]
	\arrow["{\psi|_k}", from=1-1, to=2-2]
	\arrow[dashed, from=2-1, to=2-2]
    \end{tikzcd}\]
    rendering the diagram commutative. This unique dashed arrow is precisely $\psi$ since precomposing with the injection $\Delta_k^n\times\Delta_k$ is the same as restricting along the $k$-indexed component. Thus, $\psi$ arises by universality in $\Top$, and is therefore continuous. We thus obtain that $\varphi$ is continuous, as required.

    We now show $\varphi$ is an isomorphism by constructing an explicit inverse. Define the map $\eta$ as the composite
    \[\begin{tikzcd}[row sep=0em]
	{\Delta_n} & {\displaystyle{\coprod_{[k]\in\BD}\Delta_k^n\times\Delta_k}} & {\left|\Delta^n\right|} \\
	p & {(1,p)\in\Delta_n^n\times\Delta_n} & {[(1,p)].}
	\arrow[from=1-1, to=1-2]
	\arrow[from=1-2, to=1-3]
	\arrow[maps to, from=2-1, to=2-2]
	\arrow[maps to, from=2-2, to=2-3]
    \end{tikzcd}\]
    To be fully clear, we are mapping the point $p$ to the point $p\in\Delta_n$ in the geometric $n$-simplex indexed by the identity map $1\in\Delta_n$. First, we show this map is continuous. Note that given the diagram of solid arrows
    \[\begin{tikzcd}
	& {\Delta_n^n} \\
	{\Delta_n} & {\Delta_n^n\times\Delta_n} & {\displaystyle{\coprod_{[k]\in\BD}\Delta_k^n\times\Delta_k}} & {\left|\Delta^n\right|} \\
	& {\Delta_n}
	\arrow["{x\mapsto1}", from=2-1, to=1-2]
	\arrow[dashed, from=2-1, to=2-2]
	\arrow["1"', from=2-1, to=3-2]
	\arrow[from=2-2, to=1-2]
	\arrow["{i_n}", from=2-2, to=2-3]
	\arrow[from=2-2, to=3-2]
	\arrow["\pi", from=2-3, to=2-4]
    \end{tikzcd}\]
    universality of product induces a unique dashed arrow as indicated, rendering the diagram commutative. We also clearly see that $\eta$ can be defined as the composite along the main row, where $\pi$ denotes the canonical projection and $i_n$ denotes the inclusion into the coproduct at the $n$-indexed component. Since $\pi$ and $i_n$ are continuous, we simply need to show the dashed arrow is continuous, and by universality it suffices to show the identity and the constant map at 1 are both continuous. Both are clearly continuous, giving us that $\eta$ is continuous, as required.
    
    Now we show $\eta$ is an inverse of $\varphi$. It's clearly a right inverse of $\varphi$:
    \[
    (\varphi\eta)(p)=\varphi[(1,p)]=\Delta_\bullet(1)(p)=1(p)=p,
    \]
    because $\Delta_\bullet$ is a functor and thus preserves identities. To show it is a left inverse, note that
    \[
    (\eta\varphi)[(f,p)]=\eta(\Delta_\bullet(f)(p))=[(1,\Delta_\bullet(f)(p))].
    \]
    We thus need to show that $(f,p)\sim(1,\Delta_\bullet(f)(p))$. Indeed, since $f=f^*(1)$, we have
    \[
    (f,p)=(f^*(1),p)\sim (1,\Delta_\bullet(f)(p)),
    \]
    concluding the proof.
\end{proof}

\begin{remark}\label{3.2.8}
    \Cref{3.2.7} justifies our overloaded notation for $\Delta^n$ and $\Delta^\bullet$. Because of this fact, many authors prefer to denote the standard geometric $n$-simplex by $|\Delta^n|$. We prefer this notation as well, but both choices overload notation in some way (until we prove \cref{3.2.7} at least), and we felt that for a basic introduction, our choice of notation was, for lack of a better word, less bad. However, now that we have proven \cref{3.2.7}, \textbf{we now prefer to switch to the more common notation} of denoting the standard geometric $n$-simplex by $\left|\Delta^n\right|$.
\end{remark}

\begin{remark}\label{3.2.9}
	Geometric realizations of simplicial sets are always compactly generated, and in particular are always CW complexes with one $n$-cell for each nondegenerate $n$-simplex of the simplicial set. This is extremely convenient, since the category $\mathcal{C} \mathcal{G} $ is notably much better behaved than the general category $\Top $, while still containing most notable topological spaces, making working in the category of compactly generated spaces particularly appealing.
\end{remark}

\subsection{The nerve of a category}\label{sec3.3}
Given a category $\C$, there is an extremely important construction known as the \emph{nerve} of $\C$, which encodes \emph{all} of the information of $\C$ into a simplicial set. This is an extremely important construction for a number of reasons. For example, it's proven in \cite[Prop.~1.1.2.2]{htt}\footnote{Joyal proved this first, but I prefer Lurie's treatment. However, this fact is so fundamental that it's likely a proof may date back to Boardman and Vogt, but I unfortunatey don't have access to their paper.} that all nerves are $\infty$-categories, and in particular the nerve is the canonical way of regarding categories as $\infty$-categories. Although exploring this precisely is beyond the scope of this paper, we trust that the reader understands this is an important result, even if they do not quite understand it. With this in mind, we present the definition of the nerve.

\begin{definition}\label{3.3.1}
    The nerve $N(\C)$ of a category $\C$ is the simplicial set defined as follows.
    \begin{enumerate}
        \item $n$-simplices are composable $n$-tupes of morphisms in $\C$. More precisely, the set $N(\C)_n$ of $n$-simplices is the set of all $n$-tuples $(f_0,\ldots,f_{n-1})$ of composable arrows between $(n+1)$ objects
        \[\begin{tikzcd}
            c_0 \ar["f_0",r] & c_1 \ar["f_1",r] & \cdots \ar["f_{n-2}", r] & c_{n-1} \ar["f_{n-1}", r] & c_n.
        \end{tikzcd}\]
        The set of $0$-simplices is thus simply the collection of objects $\Obj(\mathcal{C} ) $.
        \item Degeneracy maps $s_i : N(\C)_n \to N(\C)_{n+1}$ are defined by inserting the identity morphism as the $i$th arrow. Viewed as a $n$-tuple, this looks like the map
        \[
        (f_0,\ldots,f_{i-1},f_i,f_{i+1},\ldots,f_{n-1}) \mapsto (f_0,\ldots,f_{i-1},1_{c_i},f_i,f_{i+1},\ldots,f_{n-1}).
        \]
        Viewed as a diagram, this looks like taking the diagram of composable arrows
        \[\begin{tikzcd}
            c_0 \ar[r] & \cdots \ar["f_{i-1}", r] & c_{i} \ar["f_i",r] & \cdots \ar[r] & c_n
        \end{tikzcd}\]
        and mapping it to the diagram
        \[\begin{tikzcd}
            c_0 \ar[r] & \cdots \ar["f_{i-1}", r] & c_{i} \ar["1_{c_i}", r] & c_i \ar["f_i",r] & \cdots \ar[r] & c_n.
        \end{tikzcd}\]
        \item Face maps have different defintions for ``inner'' and ``outer'' faces. An inner face is a map $d_i$ where $0<i<n$, and an outer face is $d_i$ where $i=0$ or $i=n$.
        \begin{enumerate}
            \item Inner faces $d_i : N(\C)_n \to N(\C)_{n-1}$ are defined by composing the $i$th and the $(i+1)$th morphisms. Viewed as a map of $n$-tupes, this is the map
        \[
        (f_0,\ldots,f_{i-1},f_i,f_{i+1},\ldots,f_{n-1}) \mapsto (f_0,\ldots,f_{i-1},f_{i+1}\circ f_i,\ldots,f_{n-1}).
        \]
        Viewed as a diagram, this looks like taking the diagram of composable arrows
        \[\begin{tikzcd}
            c_0 \ar[r] & \cdots \ar[r] & c_{i-1} \ar["f_{i}", r] & c_{i} \ar["f_{i+1}",r] & c_{i+1} \ar[r] & \cdots \ar[r] & c_n
        \end{tikzcd}\]
        and mapping it to the diagram
        \[\begin{tikzcd}
            c_0 \ar[r] & \cdots \ar[r] & c_{i-1} \ar["f_{i+1}\circ f_{i}", r] & c_{i+1} \ar[r] & \cdots \ar[r] & c_n.
        \end{tikzcd}\]
        \item Outer faces are defined by deleting a morphism at the end of the chain. The 0th face deletes the first morphism, and the $n$th face deletes the last morphism. For example, if
        \[\begin{tikzcd}
            c_0 \ar[r] & c_1 \ar[r] & \cdots \ar[r] & c_{n-1}\ar[r] & c_n
        \end{tikzcd}\]
        is a diagram of composable arrows in $\C$, then the 0th face maps it to
        \[\begin{tikzcd}
        c_1 \ar[r] & \cdots \ar[r] & c_{n-1}\ar[r] & c_n
        \end{tikzcd}\]
        and the $n$th face maps it to
        \[\begin{tikzcd}
            c_0 \ar[r] & c_1 \ar[r] & \cdots \ar[r] & c_{n-1}.
        \end{tikzcd}\]
        \end{enumerate}
    \end{enumerate}
\end{definition}

\begin{example}\label{3.3.2}
    Consider the poset category $\mathbf{n}$. A 0-simplex $k$ in the nerve $N(\mathbf{n})$ is simply an object of $\mathbf{n}$, equivalently a non-negative integer $k<n$. A 1-simplex $f$ in the nerve is a single morphism in $\mathbf{n}$, which corresponds to a relation $h\leq k$ for $h,k\in\mathbf{n}$. Thus, 1-simplices are relations $h\leq k$ for objects $h,k$. 2-simplices are pairs of morphisms $h\xlongrightarrow{f}k\xlongrightarrow{g}\ell$, which correspond to relations $h\leq k\leq \ell$. In general, we find that $m$-simplices are relations
    \[
    h_0\leq h_1\leq\cdots\leq h_m,
    \]
    where $h_i\in\mathbf{n}$ for all $i$.

    Since $\mathbf{n}$ is well-ordered, we can say slightly more about it. Recall that \cref{3.1.6} defined \emph{degenerate} simplices as simplices which are the image of a degeneracy. In view of \cref{3.3.1}, degeneracies in the nerve simply add identity arrows. Thus any $m$-simplex with an identity arrow is automatically a degenerate simplex, since by removing the identity arrow we get an $(m-1)$-simplex, and the image of that $(m-1)$-simplex under the proper degeneracy map simply re-inserts the identity morphism back where it was. Since there are no non-identity endomorphisms in $\mathbf{n}$, note that the longest possible chain of relations in $\mathbf{n}$ without using an identity arrow is the sequence of relations
    \[
    0\leq 1\leq2\leq\cdots\leq n-2\leq n-1.
    \]
    This sequence of morphisms has length $n-1$, and it is the unique chain of arrows of length $n-1$ that can be formed without using identities. Thus, $N(\mathbf{n}$) has precisely 1 nondegenerate $(n-1)$-simplex. Moreover, no arrows longer than length $n-1$ can be made without using identities, so $N(\mathbf{n})$ has \emph{no} nondegenerate $k$-simplices for $k\geq n$.
\end{example}

\begin{example}\label{3.3.3}
    Let $G$ be a group. You have likely seen that we can view a group $G$ as a category $\mathbf{B}G$ by taking $\mathbf{B}G$ to have only one object, usually denoted $*$, by taking each morphism of $\mathbf{B}G$ to be an element of $G$, and by taking composition in $\mathbf{B}G$ to be the group operation in $G$. In particular, every morphism in $\mathbf{B}G$ has an inverse by the group axioms for $G$. We call categories where all morphisms are invertible \emph{groupoids}, as an analogue of a group. Most authors prefer to identify a group $G$ with its category, and will refer to both of them simply as $G$. Some authors prefer to write $\mathbf{B}G$ for the category, since the category can be identified with it's \emph{delooping} in $\infty$-$\G\mathrm{rpd}$. Although we prefer the former notation, we believe the latter reduces the opportunity for confusion, so we use it here for pedagogical purposes.

    We can compute the nerve of a group $G$ by considering it as a category. Since $\mathbf{B}G$ has one object, $N(\mathbf{B}G)_0=\{*\}$. The 1-simplices are the morphisms, which are precisely the elements of $G$. Thus, $N(\mathbf{B}G)_1=G$. The 2-simplices are chains of two morphisms, which are equivalently just pairs of elements in $G$, so we obtain $N(\mathbf{B}G)_2=G\times G$. This very clearly continues to give us $N(\mathbf{B}G)_n=G^n$, where $G^n$ is the $n$-repeated product. Indeed, this even generalizes to the $0$-simplices, since the empty product in $\Set$ is the terminal object, which is the singleton. Similarly to \cref{3.3.2}, we see that degeneracies insert the identity element $e\in G$, so degenerate simplices in $N(\mathbf{B}G)$ are simply tuples that contain the identity element.
\end{example}

\begin{exercise}\label{3.3.4}
    Compute the nerves of a small groupoid $\G$ and the simplex category $\BD$. Where possible, give descriptions of the degenerate simplices.
\end{exercise}

\begin{example}\label{3.3.5}
    The geometric realizations of the nerves of various categories are oftentimes fairly important objects. For example, the geometric realization of the nerve of a group is the \emph{classifying space} $BG$ of a group $G$. We will compute the very simple example $B\mathbb{Z}_2$ by following \cref{3.3.2}. We begin by constructing the coproduct
    \[
    \coprod_{[n]\in\BD} \mathbb{Z}_2^n\times\left|\Delta^n\right|.
    \]
    We will start by working only with lower dimension terms in the coproduct, to get a sense of how it will work. Since each $x\in\mathbb{Z}_2^n$ indexes a geometric $n$-simplex, we will adopt the convention of writing $|\Delta^n_{(-)}|$ to denote that the geometric $n$-simplex $|\Delta^n|$ is indexed by the $n$-simplex $(-)$. Since $|\mathbb{Z}_2|=2$, we have that the order $\left|\mathbb{Z}_2^n\right|=2^n$. We thus have
    \[
    \coprod_{[n]\in\BD} \mathbb{Z}_2^n\times|\Delta^n|\cong \left(|\Delta^0_*|\right)\amalg\left(|\Delta^1_0|\amalg|\Delta^1_1|\right)\amalg\left(|\Delta^2_{(0,0)}|\amalg|\Delta^2_{(0,1)}|\amalg|\Delta^2_{(1,0)}|\amalg|\Delta^{2}_{(1,1)}|\right)\amalg\cdots
    \]
    We begin by quotienting by degeneracies. Since $0\in\mathbb{Z}_2$ is the identity, any geometric simplex indexed by a tuple containing 0 is indexed by a degenerate simplex. For example, the simplex $|\Delta^2_{(0,0)}|$ is index by $(0,0)=s_0(0)=s_1(0)$, so $|\Delta^2_{(0,0)}|$ is identified with $|\Delta^1_0|$. This means that we should identify these geometric simplices.

    I claim we can simply drop the higher dimension simplex in the coproduct up to homeomorphism. The proof is somewhat involved. By universality of coproduct, we can simply consider the space $|\Delta_0^1|\amalg|\Delta_{(0,0)}^2|$, since a given homeomorphism, along with the identity arrows for all other maps, assemble into a homeomorphism by universality. We want to quotient this space by the relation $(s_0(x),p)\sim (x,s^0(p))$ for $x\in \mathbb{Z}_2^1$ and $p\in|\Delta_{(0,0)}^2|$. Note that the following proof is invariant under whether or not we choose the 0th or 1st degeneracy, so I chose the 0th. Our choice of $x$ is $0$, since $x$ is simply the indexing simplex for $|\Delta_0^1|$, and each point $p$ can be represented as a triple $(t_0,t_1,t_2)$ of nonnegative numbers which sum to 1. Since $s_0(0)=(0,0)$, and since $s^0(t_0,t_1,t_2)=(t_0+t_1,t_2)$, we thus have the relation
    \[
    (0,(t_0+t_1,t_2))\sim((0,0),(t_0,t_1,t_2)).
    \]
    We want to show there is a homeomorphism
    \[
    \left(|\Delta_0^1|\amalg|\Delta_{(0,0)}^2|\right)/\tildesim \cong |\Delta_0^1|.
    \]
    To prove this, define the map
    \[\begin{tikzcd}[row sep=0em]
	{    f : \left(|\Delta_0^1|\amalg|\Delta_{(0,0)}^2|\right)/\tildesim} & {|\Delta_0^1|} \\
	{[(0,(t_0,t_1))]} & {(0,t_0,t_1)}
	\arrow[from=1-1, to=1-2]
	\arrow[maps to, from=2-1, to=2-2]
\end{tikzcd}\]
    We must first show this map is well-defined. In particular, we must show each class $[(0,(t_0,t_1))]$ has precisely one representative of the given form. To do this, note that the class admits the definition as
    \[
    [(0,(t_0,t_1))] = \{(0,t_0,t_1)\}\cup\left\{(p_0,p_1,p_2)\in\Delta_{(0,0)}^2 \mid p_0+p_1=t_0,\text{ and }p_2=t_1\right\}.
    \]
    Since we are required to fix $p_0+p_1=t_0$ and $p_2=t_1$, clearly this equivalence class only has one representative of the given form, so the map is indeed well-defined. To show the map is continuous, we note that we can view the map as a map
    \[\begin{tikzcd}      
	{    |\Delta_0^1|\amalg|\Delta_{(0,0)}^2|}\ar[r] & {|\Delta_0^1|} \\
    \end{tikzcd}\]
    by noting that
    \begin{itemize}
        \item the map is the identity on the $|\Delta^1_0|$ component of the coproduct,
        \item and the map is $s^0$ on,
    \end{itemize}
    so universality of coproduct assembles our map via the unique dashed arrow in the diagram of solid arrows
    \[\begin{tikzcd}
	{|\Delta_0^1|} \\
	{|\Delta_0^1|\amalg|\Delta_{(0,0)}^2|} & {|\Delta_0^1|} \\
	{|\Delta_{(0,0)}^2|}
	\arrow[from=1-1, to=2-1]
	\arrow["1", from=1-1, to=2-2]
	\arrow[dashed, from=2-1, to=2-2]
	\arrow[from=3-1, to=2-1]
	\arrow["{s^0}"', from=3-1, to=2-2]
\end{tikzcd}\]
    rendering the diagram commutative. This map manifestly identifies elements which are related to each other under the relation $\sim$, so it factors through the quotient
    \[\begin{tikzcd}
	{|\Delta_0^1|} \\
	{|\Delta_0^1|\amalg|\Delta_{(0,0)}^2|} & {\left(|\Delta_0^1|\amalg|\Delta_{(0,0)}^2|\right)/\tildesim} & {|\Delta_0^1|} \\
	{|\Delta_{(0,0)}^2|}
	\arrow[from=1-1, to=2-1]
	\arrow[from=1-1, to=2-2]
	\arrow["1", bend left, from=1-1, to=2-3]
	\arrow[dashed, from=2-1, to=2-2]
	\arrow["f", dashed, from=2-2, to=2-3]
	\arrow[from=3-1, to=2-1]
	\arrow[from=3-1, to=2-2]
	\arrow["{s^0}"', bend right, from=3-1, to=2-3]
\end{tikzcd}\]
    Since we see that our map $f$ arises from universality conditions between continuous maps, and since $\Top$ is complete and cocomplete, necessarily $f$ is continuous.

    It thus suffices to show $f$ has a continuous inverse. The inverse is obvious, we simply define $g : (0,(t_0,t_1))\mapsto [0,(t_0,t_1)]$, and this is automatically well-defined since we showed earlier each equivalence class has precisely one representative of this form. $f$ and $g$ are manifeslty inverse functions, so we now show $g$ is continuous. Note that $g$ arises equivalently as the composite
    \[\begin{tikzcd}
        {|\Delta_0^1|}\ar["i",r] & {|\Delta^1_0|\amalg|\Delta_{(0,0)}|}\ar["\pi",r] & \displaystyle{\left(|\Delta_0^1|\amalg|\Delta_{(0,0)}^2|\right)}
    \end{tikzcd},
    \]
    where $i$ is the inclusion and $\pi$ is the natural projection onto the quotient. Since $g$ is a composite of continuous maps, it is indeed continuous.

    We have thus shown that we may freely drop degenerate simplices (almost... see \cref{3.3.6}) in the geometric realization, since that is precisely the same as quotienting by degeneracies. Thus quotienting by degeneracies is homeomorphic to the space
    \[
    \coprod_{[n]\in\BD}|\Delta^n|.
    \]
    We now need to quotient by faces. We can show that every single non-degenerate simplex is a face of a single non-degenerate simplex in precisely two ways. Let $(1,\ldots,1)_{n-1}$ be the nondegenerate $(n-1)$-simplex, and let $(1,\ldots,1)_n$ be the nondegenerate $n$-simplex. Taking an inner face map on $(1,\ldots,1)_n$ corresponds to composing two morphisms, or in other words, taking the group operation on two elements. Since $1+1=0$, this can be viewed as removing two 1s and replacing it with a 0, turning it into a degenerate simplex. However, the outer faces simply remove one of the outer 1s, and thus map $(1,\ldots,1)_n\mapsto(1,\ldots,1)_{n-1}$. Thus the only faces that are not accounted for by the degeneracy quotient are the outer faces.

    Since we only have a single copy of each geometric $n$-simplex left, from now on we will denote elements $p\in|\Delta^n|$ as $(n,p)$; in other words, the first slot in the pair is the dimensionality of the simplex the second slot lives in as a point.
    
    We begin by working inductively, and we start with the $0$-skeleton $|\Delta^0|$, which is the singleton $*$. Moving on to the 1-skeleton, we see that the inner face maps
    \[\begin{tikzcd}
	{|\Delta^0|} & {|\Delta^1|}
	\arrow["{d^0}", shift left, from=1-1, to=1-2]
	\arrow["{d^1}"', shift right, from=1-1, to=1-2]
\end{tikzcd}\]
    insert $*$ at $0$ and $1$ in the 1-simplex. Quotienting thus identifies $(1,0)\sim *\sim (1,1)$; in other words we have identified the endpoints of the 1-simplex. Since the 1-simplex is simply a line segment, this construction is homeomorphic to the unit circle $S^1$. Thus, the 1-skeleton is simply $S^1$. Now we move on to the 2-skeleton. The inner faces (subject to the previous quotient) $d^0,d^2 : S^1 \to |\Delta^2|$ insert the 1-skeleton at the $0$th and $2$nd faces. Gluing these faces together along the circle is precisely the antipodal map, producing the real projective plane $\mathbb{R}P^2$. One can induct on this construction and obtain that the $n$-skeleton is $\mathbb{R}P^n$. The geometric realization can then be seen as the colimit over the inclusions of the skeletons
    \[
    |N(\mathbf{B}G)| = \colim(X_0\hookrightarrow X_1\hookrightarrow X_2\hookrightarrow\cdots),
    \]
    which is precisely the definition of $\mathbb{R}P^{\infty}$. Thus, $B\mathbb{Z}_2=\mathbb{R}P^\infty$.
\end{example}

\begin{remark}\label{3.3.6}
	The reader might think that the section of \cref{3.3.5} stating that we may drop degenerate simplices likely generalizes, since we rarely appealed to the specific problem of \cref{3.3.5} in that section. Although our proof is sufficient for the purposes of \cref{3.3.5}, it fails in general due to transfinite considerations. We proved in \cref{3.3.5} that for any degenerate $n$-simplex, we may drop the corresponding geometric $n$-simplex. Note that this only proves that we may drop one at a time. We can induct on the argument and find that we may drop any finite number $k$ of degenerate simplices at a time, but this still fails in general since a simplicial set need not have finitely many degenerate $n$-simplices. The proof only held for \cref{3.3.5} since there are $\left| \mathbb{Z}^n_2 \right| -1=2^n-1$ degenerate $n$-simplices, which is clearly less than $\infty $, and inducting on $n$ gives us that we may drop these for each $X_n$.

	The general result is stated in \cref{3.3.7}, but the proof is notably longer. For this reason, we delegate the proof to \cref{secB} for the interested reader.
\end{remark}

\begin{proposition}\label{3.3.7}
	Let $X_\bullet$ be a simplicial set. If $X_n^{nd} $ is the set of nondegenerate $n$-simplices of $X_{\bullet} $, define the equivalence relation $\sim_f $ on $\coprod_{[n]\in\boldsymbol{\Delta} } X_n^{nd} \times \left| \Delta^n \right|$ as the relation generated by the rule $(y,\mu_*(p))\sim_f(z,\beta_*(p))$ if
	\begin{itemize}
		\item $\mu$ is a composite of face maps,
		\item $\beta$ is a composite of codegeneracy maps,
		\item $z$ is nondegenerate,
		\item $\beta^*(z)=\mu^*(y)$.
	\end{itemize}
	Let $\sim$ be the equivalence relation of \cref{3.2.2} Then there is a homeomorphism
	\[
		\left| X_{\bullet}  \right| =\left( \coprod_{[n]\in\boldsymbol{\Delta} } X_n\times \left| \Delta^n \right|  \right) /\tildesim \cong \left( \coprod_{[n]\in\boldsymbol{\Delta} } X_n^{nd} \times \left| \Delta^n \right|  \right) /\tildesim _f
	.\]
\end{proposition}
\begin{proof}
	Proven in \cref{secB}.
\end{proof}

\begin{corollary}\label{3.3.8}
    Let $X_\bullet$ be an $k$-skeletal simplicial set, as defined in \cref{3.1.10}. Then its geometric realization $|X_\bullet|$ is homeomorphic to the geometric realization of its $k$-truncation $|X_{\leq k}|$:
    \[
	    |X_{\leq k}| \coloneqq\int^{[n]\in\BD_{\leq k}} X_n\times|\Delta^n|\cong \left( \coprod_{n\leq k} X_n\times \left| \Delta^n \right|  \right) /\tildesim .
    \]
    For those unfamiliar with coends, the only change here is that the limit is taken over $\BD_{\leq k}$ as defined in \cref{3.1.10}.
\end{corollary}

\begin{exercise}\label{3.3.9}
Fun fact: we never proved the nerve was a simplicial set. Prove that the faces and degeneracies of the nerve as defined in \cref{3.3.1} satisfy the simplicial identities (\cref{3.1.7}). Afterwards, let $X$ be a set, and recall by \cref{3.1.9} that we may view $X$ as a simplicial set. We may also view it as a category by considering the discrete category structure on $X$. Compute the nerve $N(X)$, then apply \cref{3.3.8} to show $|N(X)|\cong |X|\cong X$, where the latter $X$ is viewed as a topological space with the discrete topology.
\end{exercise}

\begin{exercise}[Advanced]\label{3.3.10}
	Prove that the geometric realization of the nerve of the poset category $\mathbf{n} =[n+1]$ is the geometric realization of the standard $(n+1)$-simplicial set $|N(\mathbf{n} )|\cong \left| \Delta^{n+1}  \right| $.
\end{exercise}


\subsection{Boundaries and horns}\label{sec3.4}
We now introduce the notion of a simplicial subset, along with two important constructions known as the \emph{horn}, and the \emph{boundary/simplicial sphere}, which are of particular importance in a variety of contexts.

\begin{definition}\label{3.4.1}
    Let $X_\bullet$ be a simplicial set. A \emph{simplicial subset} of $X_\bullet$ is a simplicial set $Y_\bullet$ where
    \begin{itemize}
        \item For all $n$, the $n$-simplices of $Y_\bullet$ are a subset of those of $X_\bullet$, meaning $Y_n\subseteq X_n$;
        \item The faces and degeneracies in $Y_\bullet$ agree with those in $X_\bullet$; in particular, the faces $d_i^Y$ of $Y_\bullet$ are defined as the restrictions of the faces $d_i^X$ and degeneracies $s_i^X$ of $X_\bullet$ along the subset $Y_n\subseteq X_n$
        \[
        d_i^Y = d_i^X|_{Y_n}, \qquad s_i^Y = s_i^X|_{Y_n}.
        \]
    \end{itemize}
    We write $\Y\subseteq\X$ to denote that it is a simplicial subset.
\end{definition}

\begin{remark}\label{3.4.2}
    Unfortunately, not all collections of subsets $Y_n\subset X_n$ of a simplicial set $X_\bullet$ form a simplicial subset. Yes, one can define the faces and degeneracies for $Y_\bullet$ by restricting those of $X_\bullet$ along the relevant subsets, but \emph{the image of each restriction must be contained in the corresponding subset $Y_n$}.
    
    For example, consider the discrete simplicial set $X_\bullet$ where $X_n=\mathbb{Z}$ for all $n$, and all faces and degeneracies are the identities. We can construct subsets $Y_\bullet$ by taking $Y_n=X_n$ for all $n$ \emph{except} for $n=1$, where we take $Y_n=\mathbb{N}$. It doesn't matter whether or not we choose $0\in\mathbb{N}$ here. By \cref{3.4.1}, for $Y_\bullet$ to be a simplicial subset, the faces and degeneracies must agree with those of $X_\bullet$; in particular they are the restrictions along the subsets $Y_\bullet$. Since $Y_0=X_0$, the face $d_0 : Y_0 \to Y_1$ must be the same. Since the face is the identity, it maps each element to itself. Note that $-1\in Y_0=\mathbb{Z}$, and $d_0(-1)=-1\notin Y_1=\mathbb{N}$. Since the image is not in the subset $Y_1$, $\Y$ cannot be a simplicial subset, since the definition of $\Y$ as given does not form a valid simplicial set.
\end{remark}

\begin{remark}\label{3.4.3}
    In contrast to \cref{3.4.2}, it turns out that all one has to check to determine whether an object $Y_\bullet$ is a simplicial subset of a simplicial set $\X$ is
    \begin{itemize}
        \item $Y_n\subseteq X_n$ for all $n$;
        \item The faces and degeneracies agree with $X$;
        \item The faces and degeneracies are well-defined in the sense of \cref{3.4.2}, meaning the sets $Y_n$ contain the images of all faces and degeneracies.
    \end{itemize}
    This is fairly easily shown. The first two criteria guarentee that if $\Y$ is a simplicial set, it will be a simplicial subset of $\X$, so we need only check that it is a simplicial set. Since the faces and degeneracies are all well-defined, \cref{3.1.7} tells us that we need only check the faces and degeneracies satisfy the simplicial identities. This is manifestly true: the faces and degeneracies in $\Y$ are simply restrictions of those on $\X$, and since $\X$ is a simplicial set, the faces and degeneracies satsify the simplicial identities on the whole of $\X$. This means they must satisfy them for every element of each $X_n$, and thus every subset of each $X_n$.

    In particular, we see that to specify a simplicial subset $\Y$ of $\X$, one needs only specify the sets $Y_n$, and then show these sets contain the images of the restricted faces and degeneracies, since the definitions of the faces and degeneracies are forced by \cref{3.4.1}.
\end{remark}

\begin{definition}\label{3.4.4}
    Let $\alpha : \X\to\Y$ be a morphism of simplicial sets. The \emph{image} of $\alpha$, denoted $\im\alpha_\bullet$, is defined as the simplicial subset of $\Y$ given by
    \[
    (\im\alpha)_n=\im(\alpha_n : X_n\to Y_n),
    \]
    where $\alpha_n$ is the component of $\alpha$ at $n$. We make use of the convenient notation and write $\im\alpha_n$ to denote the set of $n$-simplices of $\im\alpha$.
\end{definition}

\begin{claim}\label{3.4.5}
    With notation as in \cref{3.4.4}, the image is indeed a simplicial subset of $\Y$.
\end{claim}
\begin{proof}
    By \cref{3.4.3}, we need only show the images of each face and degeneracy land within the specified sets. We start by showing this for degeneracies. Let $y\in\im\alpha_n$, and we want to show $s_i^{\im\alpha}(y)\in\im\alpha_{n+1}$. Recall by \cref{3.1.5} that morphisms of simplicial sets are natural tranformations, meaning they commute with degeneracies as indicated in the diagram
    \[\begin{tikzcd}
	{X_n} & {Y_n} \\
	{X_{n+1}} & {Y_{n+1}}
	\arrow["{\alpha_n}", from=1-1, to=1-2]
	\arrow["{s_i^X}"', from=1-1, to=2-1]
	\arrow["{s_i^Y}", from=1-2, to=2-2]
	\arrow["{\alpha_{n+1}}"', from=2-1, to=2-2]
    \end{tikzcd}\]
    By definition, for each $y\in\im\alpha_n$, there exists at least one $x\in X_n$ such that $y=\alpha_n(x)$. Since the degeneracy in $\im\alpha_\bullet$ is the same as the degeneracy for $\Y$, we have
    \[
    s^{\im\alpha}_i(y)=s^Y_i(y)=\left(s^Y_i\alpha_n\right)(x)=\left(\alpha_{n+1}s_i^X\right)(x),
    \]
    where the final step follows by commutativity of the diagram. By definition, this is in the image of $\alpha_{n+1}$, completing the proof for degeneracies. The proof for faces is the same, and is done by constructing the corresponding naturality diagram for faces instead of degeneracies.
\end{proof}

\begin{definition}\label{3.4.6}
    Let $\X,\Y$ be simplicial sets, and suppose that for each $n$, the faces and degeneraces of $\X$ and $\Y$ agree on the intersection $X_n\cap Y_n$. In symbols,
    \[
    s^Y_i|_{X_n\cap Y_n}=s_i^X|_{X_n\cap Y_n},\qquad d^Y_i|_{X_n\cap Y_n}=d_i^X|_{X_n\cap Y_n}.
    \]
    Then the \emph{union} $\X\cup\Y$ of $\X$ and $\Y$ is defined as the levelwise union
    \[
    (X\cup Y)_n=X_n\cup Y_n,
    \]
    and the faces and degeneracies are defined by
    \[
    s_i(x) = \begin{cases}
        s_i^X(x),&x\in X\\
        s_i^Y(x),&x\in Y
    \end{cases},\qquad d_i(x) = \begin{cases}
        d_i^X(x),&x\in X\\
        d_i^Y(x),&x\in Y
    \end{cases}.
    \]
\end{definition}

\begin{claim}\label{3.4.7}
    With notation as in \cref{3.4.6}, $\X\cup\Y$ is a simplicial set. In particular, $\X,\Y$ are simplicial subsets of their union.
\end{claim}
\begin{proof}
    The fact that $\X,\Y\subseteq\X\cup\Y$ follows manifestly from the definition, provided we indeed show $\X\cup\Y$ is a simplicial set, which is not terribly difficult. As usual, we appeal to \cref{3.1.7} and show that the faces and degeneracies are well defined, and satisfy the simplicial identities. Since $\X$ and $\Y$ agree on their intersection, for any $x\in X_n\cap Y_n$, we must have $d_i^X(x)=d_i^Y(x)$, and the same for degeneracies. Thus, the maps are well-defined. To show they satisfy the simplicial identities, note that for all $x\in X_n\cup Y_n$, we have $x\in X_n$ or $x\in Y_n$. Without loss of generality, let $x\in X_n$. Since $\X$ is a simplicial set, the simplicial operators satisfy the simplicial identities for $x$. Extending this argument to $\Y$, we find that the simplicial operators satisfy the simplicial identities for all $x\in X_n\cup Y_n$, and thus $\X\cup\Y$ is indeed a simplicial set.
\end{proof}

Previously in \cref{2.1.6}, we defined the \emph{face} of the standard $n$-simplex, and we were even able to encode this information into the category $\BD$ in \cref{2.4.4}. As explained in \cref{3.2.5}, the combinatorial and geometric information encoded by simplicial sets is in many ways more general than that of $\BD$, and as seen in \cref{3.2.7}, the standard $n$-simplicial sets encode the information of the standard geometric $n$-simplex. We should thus think that there should be some way to define the faces of the standard $n$-simplicial set. Indeed there is such a way to do so, and the definition of the face of the standard $n$-simplicial set will give rise to much more important constructions.

\begin{definition}\label{3.4.8}
    Let $d^i$ be the image of the coface map $[n-1]\to[n]$ under the Yoneda embedding. More precisely, let $d^i : \Delta^{n-1}\to\Delta^n$ denote precomposition by the coface map $d^i:[n-1]\to[n]$.
    \begin{enumerate}
        \item The $i$th face $\partial_i\Delta^n$ of the simplical set $\Delta^n$ is the simplicial subset $\im d^i_\bullet$.
        \item The \emph{boundary} $\partial\Delta^n$ of the standard $n$-simplicial set $\Delta^n$, also called the \emph{simplicial $n$-sphere}, is the simplicial subset
        \[
        \partial\Delta^n \coloneqq\bigcup_{i\in[n]} \partial_i\Delta^n=\bigcup_{i\in[n]}\im d^i_\bullet \subseteq\Delta^n.
        \]
        \item The \emph{$k$th horn} $\Lambda_n^k$ of the standard $n$-simplicial set is the simplicial subset
        \[
        \Lambda_n^k\coloneqq\bigcup_{i\in[n],i\neq k}\partial_i\Delta^n=\bigcup_{i\in[n],i\neq k}\im d^i_\bullet \subseteq\partial\Delta^n\subseteq\Delta^n.
        \]
    \end{enumerate}
\end{definition}

\begin{remark}\label{3.4.9}
    Because boundaries and horns are defined as unions over simplicial subsets of $\Delta^n$, their faces and degeneracies automatically agree with $\Delta^n$, and in particular agree with each other, making the union of simplicial sets well-defined.
\end{remark}

\begin{intuition}\label{3.4.10}
    The geometric intuition for faces has already been discussed. Our definition of faces implies that the coface maps $d^i : \Delta^{n-1}\to\Delta^n$ do the ``simplicial set equivalent'' of inserting the standard $(n-1)$-simplicial set at the ``$i$th face of the standard $n$-simplicial set''. It remains to be shown of course that the geometric realization of $\im d^i_\bullet$ is the $i$th face of $|\Delta^n|$, but this will be shown soon.

    For now, working with the understanding that $\im d^i_\bullet$ represents the face of the standard $n$-simplicial set, we define the \emph{boundary} of the standard $n$-simplicial set as the union of all of its faces. We expect a union of simplicial sets to correspond to a union of geometric objects, so working under this intuition, we would imagine the geometric realization of the boundary $\partial\Delta^n$ to be the physical boundary of the standard $n$-simplex. Once again, we will show this to be true shortly.

    The only definition that may not be immediately clear is that of a horn. The definition of the $k$th horn $\Lambda_n^k$ is the union of all the faces of $\Delta^n$ \emph{except} the $k$th one. For example, consider $\Lambda_2^2$. Geometrically, this corresponds to taking the standard $2$-simplex, deleting the interior so that we are only left with its faces (which constitute the boundary), and then deleting the $2$nd face. \Cref{fig:6} illustrates this geometrically.
    \begin{figure}[ht]
        \centering
        \begin{tikzpicture}[scale=5]
    % Set the viewing angle for the 3D plot
    \tdplotsetmaincoords{70}{110}

    % Establish the 3D coordinate system
    \begin{scope}[tdplot_main_coords]
    
        % Draw the axes for context
        \draw[-latex] (0,0,0) -- (1.2,0,0) node[anchor=north east]{$x$};
        \draw[-latex] (0,0,0) -- (0,1.2,0) node[anchor=north west]{$y$};
        \draw[-latex] (0,0,0) -- (0,0,1.2) node[anchor=south]{$z$};

        % Define the vertices of the standard 2-simplex
        \coordinate (v0) at (1,0,0);
        \coordinate (v1) at (0,1,0);
        \coordinate (v2) at (0,0,1);

        \draw[thick, blue] 
            (v1) -- (v2);
        \draw[thick, blue] 
            (v0) -- (v2);
        \draw[thin, dashed, gray] 
            (v0) -- (v1);

        % Label the vertices (0-faces)
        \node[at=(v0), anchor=south east, inner sep=2pt] {$v_0$};
        \node[at=(v1), anchor=south west, inner sep=2pt] {$v_1$};
        \node[at=(v2), anchor=south west, inner sep=5pt] {$v_2$};

        \node[at=($(v1)!0.5!(v2)$), anchor=south west, inner sep=3pt, text=blue] {$f_0$};
        \node[at=($(v0)!0.5!(v2)$), anchor=south east, inner sep=3pt, text=blue] {$f_1$};
        \node[at=($(v0)!0.5!(v1)$), anchor=north, inner sep=3pt, text=blue] {$f_2$ - Deleted face};

    \end{scope}
\end{tikzpicture}
        \caption{The second horn $|\Lambda_2^2|$ of $|\Delta^2|$}
        \label{fig:6}
    \end{figure}
    The remaining two faces, $f_0$ and $f_1$, assemble into a sort of ``$\wedge$'' shape, which looks somewhat like a two-dimensional horn. Hence the namesake. It remains to be shown of course that the geometric realizations align with these objects. We can in fact prove stronger results. In \cref{sec4.1}, we will show that geometric realization is a left adjoint functor, and thus preserves colimits. It thus suffices to show that unions and images can be viewed as colimits, and our intuition here will follow as \cref{3.4.13}.
\end{intuition}


\begin{proposition}\label{3.4.11}
    Let $F : \J\to\D^\C$ be a diagram with $\D$ complete, and let $\mathrm{ev}_c:\D^\C\to\D$ be the evaluation functor $F\mapsto F(c)$, $\alpha\mapsto\alpha_c$. Then $\lim F$, if it exists, is given by $(\lim F)(c)=\lim (\mathrm{ev}_c\circ F)$.
\end{proposition}
\begin{proof}
    See \cite[\S5.3~Thm.~1]{ml}.
\end{proof}

\begin{corollary}\label{3.4.12}
    Images and unions of simplicial sets, as defined in \cref{3.4.1} and \cref{3.4.6}, are colimits in $\sSet$. In particular, if $\X,\Y\in\sSet$ and $\alpha : \X\to\Y$ is a morphism of simplicial sets, then
    \begin{enumerate}
        \item $\left|\im(\alpha : \X\to\Y)_\bullet\right|\cong\im(\alpha : |\X|\to|\Y|)$, and
        \item $\left|\X\cup\Y\right|\cong|\X|\cup|\Y|$.
    \end{enumerate}
\end{corollary}
\begin{proof}
    Recall that $\sSet=\Set^{\BD\opp}$ is a functor category, and since $\Set$ is complete and cocomplete, \cref{3.4.11} and its dual apply. In particular, any diagram $F$ in $\sSet$ must satisfy $(\lim F)_n=\lim (\mathrm{ev}_n\circ F)$. This tells us that limit in $\sSet$ are computed levelwise, meaning in particular
    \begin{enumerate}
        \item Let $\alpha : \X\to\Y$ be a morphism of simplicial sets. The \emph{limit} definition of the image $\im\alpha_\bullet$ (not \cref{3.4.4}) must satisfy $(\im\alpha)_n=\im(\alpha_n)$, where the second is a limit in $\Set$. Note that this is precisely \cref{3.4.4}, and thus we have that our image in $\sSet$ can be regarded as a limit. Since $\Set$ admits canonical decompositions for morphisms, there is an isomorphism $\im(\alpha_n)\cong\coim(\alpha_n)$. The coimage is a colimit, and thus we may regard the image in $\sSet$ as isomorphic to a colimit.
        \item Let $\X,\Y$ be simplicial sets whose simplicial operators agree on their intersection. Because their simplicial operators agree, it is easy to see that $\X\cap\Y$ is a subsimplicial set of both $\X$ and $\Y$, with the intersection given by levelwise intersection. Note that the pushout $Z_\bullet$ of $\X$ and $\Y$ along the canonical inclusions $\X\cap\Y\hookrightarrow \X,\Y$ can be computed levelwise, as indicated in the diagram
        \[\begin{tikzcd}
	{(X\cap Y)_n} & {X_n} \\
	{Y_n} & {Z_n}
	\arrow[hook', from=1-1, to=1-2]
	\arrow[hook, from=1-1, to=2-1]
	\arrow["\lrcorner"{anchor=center, pos=0.125}, draw=none, from=1-1, to=2-2]
	\arrow[from=1-2, to=2-2]
	\arrow[from=2-1, to=2-2]
    \end{tikzcd}\]
    in $\Set$. It is well known that this pushout can be given by $Z_n=X_n\cup Y_n$, which is precisely the definition of the $n$-simplices of the union $(X\cup Y)_n$. Thus, we may regard the union as the pushout, which is in particular a colimit.
    \end{enumerate}
    Recall that geometric realiation $|-|$ is left adjoint, and thus preserves colimits. This implies that $|\X\cup\Y|\cong|\X|\cup|\Y|$ and that $|\coim\alpha|=\coim|\alpha|$. Note that $\Top$ also admits canonical decompositions of morphisms, so we obtain $\coim|\alpha|\cong\im|\alpha|$, giving a string of isomorphism
    \[
    |\im\alpha|\cong|\coim\alpha|\cong\coim|\alpha|\cong\im|\alpha|.
    \]
\end{proof}

\begin{corollary}\label{3.4.13}
    The following statements hold.
\begin{enumerate}
    \item The geometric realization of the $i$th face is the $i$th face of the geometric $n$-simplex:
    \[
    \left|\partial_i\Delta^n\right|\cong d^i(|\Delta^{n-1}|).
    \]
    \item The geometric realization of the boundary is the boundary of the geometric $n$-simplex:
    \[
    \left|\partial\Delta^n\right|\cong\partial\left|\Delta^n\right|.
    \]
    \item The geometric realization of the $k$th horn is the $k$th horn (defined as the union of all faces other than $k$) of the geometric $n$-simplex:
    \[
    \left|\Lambda_n^k\right|\cong\bigcup_{i\in[n],i\neq k}\im (d^i :\left|\Delta^{n-1}\right|\to\left|\Delta^n\right|).
    \]
\end{enumerate}
\end{corollary}
\begin{proof}
	Indeed, these are all special cases of $\cref{3.4.12}$, and it suffices to show $|d^i|=d^i : |\Delta^{n-1}|\to|\Delta^n|$. This is shown in \cref{sec3.5}.
\end{proof}

\begin{notation}\label{3.4.14}
    As is likely expected, we represent the geometric $k$th horn of the standard $n$-simplex as $|\Lambda_n^k|$. We also adopt the convention of denoting the $i$th face of the geometric $n$-simplex as $\partial_i|\Delta^n|$.
\end{notation}

\begin{remark}\label{3.4.15}
    It should be fairly clear to anyone who has taken topology that the boundary of the triangle $\partial|\Delta^2|$ is homeomorphic to the 1-sphere $S^1$. Similarly, it's clear that the boundary of the tetrahedron $\partial|\Delta^3|$ is homeomorphic to the 2-sphere $S^2$. Indeed this generalizes: for $n>1$, there is a homeomorphism $\partial|\Delta^n|\cong S^{n-1}$. For this reason, many authors call $\partial\Delta^n$ the simplicial $n$-sphere.
\end{remark}

We close this chapter by establishing a collection of important facts about horns. These facts are extremely useful in establishing basic results for \emph{Kan complexes} and \emph{quasi-categories}, which are explored a small amount in \cref{sec4.4}, and play a vital role in homotopy theory and higher category theory.

\begin{definition}\label{3.4.16}
	Let $X_{\bullet} $ be a simplicial set. A \emph{horn} of $X_{\bullet} $ is a morphism of simplicial sets $\Lambda_i^n \to X_{\bullet} $. A \emph{face} of $X_{\bullet} $ is a morphism of simplicial sets $\partial _i\Delta^n \to K$.
\end{definition}

\begin{lemma}\label{3.4.17}
	Let $X_{\bullet} $ be a simplicial set. Then a morphism $\sigma : \Delta^n \to X_{\bullet} $ of simplicial sets is simply a choice of an $n$-simplex $\tau\in X_n$, and mapping $\sigma(1_{[n]} )=\tau$.
\end{lemma}
\begin{proof}
	Indeed, this is precisely the Yoneda lemma applied to simplicial sets.
\end{proof}

\begin{remark}\label{3.4.18}
	Since the epi-mono factorization in $\boldsymbol{\Delta} $ given by \cref{2.3.4} is unique, and since elements of $\partial _i\Delta^n$ are simply given by postcomposing $\Delta^{n-1} $ by a coface, uniqueness gives an isomorphism $\partial_i\Delta^n\cong\Delta^{n-1} $ given by simply removing the coface from the composite. Applying \cref{3.4.17}, this gives us the fact that a morphism of simplicial sets $\partial _i\Delta^n \to X_{\bullet} $ is a choice of an $(n-1)$-simplex $\tau\in X_{n-1} $, and mapping $d^i \mapsto \tau$.
\end{remark}

\begin{terminology}
	In view of \cref{3.4.16,3.4.17,3.4.18}, we will often refer to morphisms $\Delta^n \to X_{\bullet} $ as $n$-simplices of $X_{\bullet} $, and will refer to morphisms $d^j\mapsto \tau\in X_n$ as \emph{faces} of $X_{\bullet} $.
\end{terminology}


\begin{proposition}\label{3.4.20}
	A horn $\Lambda_i^n \to X_{\bullet} $ in a simplicial set $X_{\bullet} $ is determined by an $(n-1)$-simplex $\tau_j$ for each $j\neq i,j\in[n]$ such that the compatibility condition $d_{j} (\tau_k)=d_{k-1} (\tau_j)$ for $j<k$ holds. The horn is then given by $d^j \mapsto \tau_j$.
\end{proposition}
\begin{proof}
	The proof is quite lengthy, so we offer only a sketch. In the proof of \cref{3.4.12}, we observed that unions in $\sSet $ are unions in the category-theoretic sense, so one constructs the relevant pullback diagram, which (by \cref{3.4.18}) simplifies the statement to ``the diagram commutes if and only if the compatibility condition holds''. If the diagram commutes, then the condition follows very easily by simplicial identities. For the converse direction, it suffices to show each face $d^j \mapsto \tau_j$ agrees on the intersection
	\[
		K_{\bullet} \coloneqq \bigcap_{j\in[n],j\neq i} \partial_j\Delta^n
	.\]
	This is shown by proving that $K_{\bullet} $ has $i$ as a unique vertex, and that it has $s_m s_{m-1} \ldots s_0(i)$ as a unique $m$-simplex. We then show the faces agree on $i$ by applying the compatibility condition on the $j$th and $(j+1)$th faces for $i<j$ and $j<i-1$; we then apply the condition for $i-1$ and $i+1$, showing any two faces agree by transitivity. By finiteness of $[n]$, induction gives that all faces agree on $i$, so it suffices to show they agree on degeneracies, which is immediate from the fact that the faces are morphisms of simplicial sets.
\end{proof}

\begin{definition}\label{3.4.21}
	Let $\sigma_0: \Lambda_i^n\to X_{\bullet} $ be a horn in $X_{\bullet} $. A \emph{filler} for the horn $\sigma_0$ is an $n$-simplex $\sigma : \Delta^n \to X_{\bullet} $ which agrees with $\sigma_0$ on the horn. More precisely,
	\[
		\sigma|_{\Lambda_i^n} =\sigma_0
	.\]
	This is often also expressed as the existence of the (not necessarily unique) dotted arrow as indicated
	\[\begin{tikzcd}[row sep = 2.5em, column sep = 2.5em]
	\Lambda_i^n \ar[" \sigma_0 ", r] \ar[hook, d]  & X_{\bullet}  \\
	\Delta^n \ar[" \sigma "', dotted, ur]  & 
	\end{tikzcd}\]
	rendering the diagram commutative.
\end{definition}


\begin{corollary}\label{3.4.22}
	Let $X_{\bullet} $ be a simplicial set. A filler for a horn $\sigma_0:\Lambda_i^n\to X_{\bullet} $ is an $n$-simplex $\sigma : \Delta^n \to X_{\bullet} $, $1_{[n]} \mapsto \tau\in X_n$, such that for each $j\neq i$ and each face $d^j : [n-1] \to [n]$,
			\[
				\sigma_0(d^j)=\sigma(d^j)
			.\]	
\end{corollary}
\begin{proof}
	The restriction $\sigma|_{\Lambda_i^n}$ gives a horn in $K$, and by \cref{3.4.20} this amounts to a choice $\tau_j$ for each $j\neq i$ satisfying the compatibiltiy condition, and mapping $d^j \mapsto \tau_j$. Also by \cref{3.4.20}, this horn agrees on (non-$i$) faces with $\sigma_0$ precisely when $\sigma|_{\Lambda_i^n} = \sigma_0$.
\end{proof}

\begin{exercise}[Challenge]\label{3.4.23}
	Technically, the reader now has the necessary knowledge to complete the following proof. This proof is fairly nontrivial, and the reader will likely struggle, especially if this is their first exposure to simplicial sets. We present the problem for the ambitious reader, but we do not expect the reader to solve it. A proof can be found in \cite[Prop.~1.1.2.2]{htt}. Let $X_{\bullet} $ be a simplicial set. Prove the following are equivalent:
	\begin{enumerate}
		\item There is a small category $\mathcal{C} $ and an isomorphism $N(\mathcal{C} )\cong X_{\bullet} $.
		\item Every \emph{inner horn} admits a \emph{unique} filler. More precisely, for all $0<i<n$, for all diagrams of solid arrows
			\[\begin{tikzcd}[row sep = 2.5em, column sep = 2.5em]
			\Lambda_i^n \ar[r] \ar[hook, d]  & X_{\bullet}  \\
			\Delta^n \ar[dashed, ur]  & 
			\end{tikzcd}\]
			there exists a \emph{unique} dashed arrow as indicated, rendering the diagram commutative.
	\end{enumerate}
\end{exercise}


